% Generated draft: requires mathematical and source review.
\documentclass[11pt]{article}
\usepackage[margin=25mm]{geometry}
\usepackage{fontspec}
\setmainfont{Latin Modern Roman}
\usepackage{amsmath,amssymb,mathtools,unicode-math}
\setmathfont{Latin Modern Math}
\usepackage{enumitem,needspace,xurl,hyperref}
\setlength{\parindent}{0pt}
\setlength{\parskip}{5pt}
\setlength{\emergencystretch}{3em}
\newcommand{\N}{\mathbb{N}}
\newcommand{\Z}{\mathbb{Z}}
\newcommand{\Q}{\mathbb{Q}}
\newcommand{\R}{\mathbb{R}}
\newcommand{\C}{\mathbb{C}}
\newcommand{\F}{\mathbb{F}}
\newcommand{\A}{\mathbb{A}}
\newcommand{\PP}{\mathbb{P}}
\newcommand{\E}{\mathbb{E}}
\newcommand{\eps}{\varepsilon}
\newcommand{\dd}{\,\mathrm{d}}
\begin{document}
\section*{UnsolvedMath ID: 701183}
\textbf{Source problem code:} catalogue entry
\subsection{Definitions and mathematical statement}
Let \(p\) be a prime, let \(\Omega\) be a nonempty set, and let \(G\leq\operatorname{Sym}(\Omega)\) be a finite \(p\)-group acting transitively on \(\Omega\). Thus the permutation action is faithful, and transitivity implies that \(\Omega\) is finite. An element \(g\in G\) is fixed-point-free if \(g(\omega)\ne\omega\) for every \(\omega\in\Omega\). Write \(|g|\) for its order and define the exponent by
\[
\exp(G)=\operatorname{lcm}\{\,|g|:g\in G\,\}.
\]
Assume that
\[
\forall g\in G,\qquad
\bigl(\forall\omega\in\Omega,\ g(\omega)\ne\omega\bigr)
\Longrightarrow |g|=p.
\]
\textbf{Question (source formulation).} Must \(\exp(G)=p\)?

\textbf{Unresolved boundary qualification.} The source does not explicitly exclude the trivial group acting on a singleton. In that case the fixed-point-free hypothesis is vacuous and the exponent is \(1\). Consequently, precision requires clarification whether nontriviality is implicit or whether the intended conclusion is instead \(\exp(G)\mid p\), equivalently \(g^p=1\) for every \(g\in G\). For nontrivial finite \(p\)-groups these conclusions coincide. No repair to the source wording is silently imposed here.

The supplied source states that the intended assertion is true for \(p=2,3\), and that under the stated hypothesis the exponent is bounded by a function of \(p\) alone. These are reported background claims, not independently verified results in this draft.
\subsection{Short English statement}
The question asks whether a finite transitive permutation \(p\)-group whose fixed-point-free elements all have order \(p\) must itself have exponent \(p\). The degree-one boundary case requires clarification. The registry labels the problem open; this draft does not independently verify that status.
\subsection{Sources}
\begin{itemize}
\item Local UnsolvedMath registry, permanent ID \(701183\), source code \texttt{catalogue entry}, title \emph{Kourovka Notebook Problem 15.95}; supplied upstream revision \texttt{d371ee917ac9647fbb06cd6fc529deec6f3dda82}.
\item \emph{The Kourovka Notebook (Unsolved Problems in Group Theory)}, original problem 15.95, 15th issue (2002); the supplied background also identifies the 21st issue edition (2026), arXiv:1401.0300v48. The collection description names E. I. Khukhro and V. D. Mazurov as editors.
\item The registry and background attribute the proposal to E. I. Khukhro, with proposed year 2002. The supplied statement separately begins with the attribution ``A. Mann, Ch. Praeger''; both attributions are retained without assigning an unprovided relationship between them.
\item A. Hassani, M. Khayaty, E. I. Khukhro, C. E. Praeger, \emph{J. Algebra} \textbf{214}, no. 1 (1999), 317--337. The source cites this reference for the assertion that the exponent is bounded in terms of \(p\). No article title or link is supplied.
\end{itemize}
\end{document}
